\documentclass[11pt]{article}

\usepackage[T1]{fontenc}
\usepackage{lmodern}
\usepackage{amsmath,amssymb,amsthm,mathtools}
\usepackage{booktabs,longtable,array}
\usepackage{enumitem}
\usepackage{geometry}
\usepackage{microtype}
\usepackage{xcolor}
\usepackage{xurl}
\usepackage{tabularx}
\usepackage{hyperref}
\usepackage{pdflscape}
\usepackage{listings}

\geometry{margin=1in}
\allowdisplaybreaks
\hypersetup{
  colorlinks=true,
  linkcolor=blue!45!black,
  citecolor=blue!45!black,
  urlcolor=blue!45!black,
  pdftitle={The Steiner--Wiener 4-index inverse problem for finite trees},
  pdfauthor={[AUTHOR NAME]}
}
\setlist{nosep}
\setlength{\emergencystretch}{3em}

\newtheorem{theorem}{Theorem}[section]
\newtheorem{lemma}[theorem]{Lemma}
\newtheorem{proposition}[theorem]{Proposition}
\newtheorem{corollary}[theorem]{Corollary}
\theoremstyle{definition}
\newtheorem{definition}[theorem]{Definition}
\newtheorem{certificate}[theorem]{Computer-assisted certificate}
\theoremstyle{remark}
\newtheorem{remark}[theorem]{Remark}

\newcommand{\SW}{\operatorname{SW}_4}
\newcommand{\Tfour}{\mathcal T_4}
\newcommand{\N}{\mathbb N}
\newcommand{\C}[2]{\binom{#1}{#2}}
\newcommand{\set}[1]{\left\{#1\right\}}
\newcommand{\interval}[2]{\left[#1,#2\right]}
\newcommand{\Mink}{\mathbin{+}}
\newcommand{\code}[1]{\nolinkurl{#1}}

\title{The Steiner--Wiener $4$-Index Inverse Problem for Finite Trees:\\
A Complete Computer-Assisted Solution}
\author{\texttt{[AUTHOR NAME]}\\
\texttt{[AFFILIATION]}\\
\texttt{[EMAIL]}\\
ORCID: \texttt{[ORCID]}}
\date{}

\begin{document}
\maketitle

\begin{abstract}
For a finite tree $T$, the Steiner--Wiener $4$-index is the sum, over all
four-element vertex sets, of the number of edges in the minimal subtree
spanning the set.  We determine exactly which positive integers fail to occur
as such an index.  The proof combines an exact rooted-forest dynamic program
covering all trees in the exceptional range, exact caterpillar subset-sum and
compressed-interval certificates for a finite bridge, and a universal
Thue--Morse caterpillar construction whose fixed-order intervals overlap for
all sufficiently large even orders.  Every finite computation uses exact
integer arithmetic and is accompanied by reproducible source, exact output,
and independent checks.  With $\N=\{1,2,3,\ldots\}$, the final result is
\[
  \left|\N\setminus\Tfour\right|=14{,}099,
  \qquad
  \max(\N\setminus\Tfour)=80{,}163.
\]
\end{abstract}

\noindent\textbf{Keywords.}
Steiner distance; Steiner--Wiener index; tree invariant; inverse problem;
caterpillar; exact certificate; Thue--Morse partition.

\noindent\textbf{2020 Mathematics Subject Classification.}
05C05, 05C12, 05C35.

\section{Introduction and main theorem}

For a finite tree $T$ and a nonempty set $S\subseteq V(T)$, let $d_T(S)$
denote the number of edges in the unique minimal connected subgraph of $T$
containing $S$.  The Steiner--Wiener $4$-index is
\begin{equation}\label{eq:def-sw}
  \SW(T)=
  \sum_{\substack{S\subseteq V(T)\\ |S|=4}}d_T(S).
\end{equation}
We write
\[
  \Tfour=\set{\SW(T):T\text{ is a finite tree}},
  \qquad
  \N=\{1,2,3,\ldots\}.
\]
The two inverse parameters are
\[
  e(4)=\left|\N\setminus\Tfour\right|,
  \qquad
  N_0(4)=\max\left(\N\setminus\Tfour\right),
\]
where the theorem below proves in particular that the complement is finite and
nonempty.

\begin{theorem}[Main theorem]\label{thm:main}
Exactly $14{,}099$ positive integers are not equal to $\SW(T)$ for any finite
tree $T$, and the largest such integer is $80{,}163$.  Equivalently,
\[
  \boxed{e(4)=14{,}099},
  \qquad
  \boxed{N_0(4)=80{,}163}.
\]
\end{theorem}

The proof has three logically distinct parts.
\begin{enumerate}[label=(\roman*)]
\item An exact dynamic program enumerates the attainable values of
      \emph{all} trees in the range in which exceptions can occur.  It proves
      that exactly $14{,}099$ values at most $80{,}163$ are absent, that
      $80{,}163$ is absent, and that every value from $80{,}164$ through
      $178{,}857$ occurs.
\item Exact caterpillar certificates and compressed interval unions prove
      \[
        \interval{178{,}858}{554{,}860{,}689{,}583}
        \subseteq\Tfour.
      \]
\item A self-contained caterpillar construction gives
      \[
        \interval{272{,}308{,}250{,}181}{\infty}\subseteq\Tfour.
      \]
      This overlaps the finite bridge.
\end{enumerate}
The universal arguments are proved symbolically.  Statements restricted to
finite parameter ranges are explicitly identified as exact computer-assisted
certificates.

\section{Edge decomposition and caterpillar coordinates}

We use the convention $\binom{m}{r}=0$ for integers $m<r$.

\subsection{The edge-contribution formula}

\begin{lemma}[Edge decomposition]\label{lem:edge}
Let $T$ be a tree of order $n$.  If deleting an edge $e$ produces components
of orders $s_e$ and $n-s_e$, then
\begin{equation}\label{eq:edge-weight}
  \SW(T)=\sum_{e\in E(T)}w_n(s_e),
  \qquad
  w_n(s)=\C{n}{4}-\C{s}{4}-\C{n-s}{4}.
\end{equation}
\end{lemma}

\begin{proof}
For a four-set $S$, the edge $e$ belongs to the minimal subtree spanning $S$
if and only if $S$ is not contained entirely in either component of $T-e$.
Among all $\binom n4$ four-sets, exactly $\binom{s_e}{4}$ lie in one
component and $\binom{n-s_e}{4}$ lie in the other.  Thus $e$ is counted
$w_n(s_e)$ times in \eqref{eq:def-sw}.  Summing the edge indicators first and
then the four-sets proves \eqref{eq:edge-weight}.
\end{proof}

The first difference is
\begin{equation}\label{eq:w-difference}
  w_n(s+1)-w_n(s)
  =\C{n-s-1}{3}-\C{s}{3}.
\end{equation}
It is nonnegative for $1\le s<n/2$.  Choosing $s_e$ as the smaller component
order therefore gives $w_n(s_e)\ge w_n(1)=\binom{n-1}{3}$.

\begin{corollary}[Minimum at a fixed order]\label{cor:star-min}
The minimum Steiner--Wiener $4$-index among trees of order $n\ge4$ is
\begin{equation}\label{eq:G}
  G(n)=(n-1)\C{n-1}{3}.
\end{equation}
For $n\ge5$, equality holds only for the star.  At $n=4$, every tree has
index $3=G(4)$.
\end{corollary}

\begin{proof}
Every $n$-vertex tree has $n-1$ edges, so the preceding inequality gives
$\SW(T)\ge(n-1)w_n(1)=G(n)$.  A star attains equality.  For $n\ge5$ the
inequality in \eqref{eq:w-difference} is strict at every $1\le s<n/2$ once
$s>1$; hence equality forces every edge to have a one-vertex component, which
characterizes the star.  The order-four statement follows directly from
\eqref{eq:edge-weight}.
\end{proof}

\subsection{Caterpillar cut sets}

The finite bridge and the infinite construction use a convenient exact
parameterization of caterpillars.

\begin{lemma}[Caterpillar cut-set realization]\label{lem:caterpillar-realization}
Let $n\ge4$ and let
$A=\{a_1<\cdots<a_t\}\subseteq\{2,\ldots,n-2\}$.  There is an $n$-vertex
caterpillar whose non-leaf edge cuts have component orders
$a_1,\ldots,a_t$.  Its index is
\begin{equation}\label{eq:caterpillar-formula}
  \SW(T_A)=G(n)+\sum_{a\in A}\Delta_n(a),
  \qquad
  \Delta_n(a)=\C{n-1}{4}-\C{a}{4}-\C{n-a}{4}.
\end{equation}
The same formula with an empty sum is realized by the star when $A=\varnothing$.
\end{lemma}

\begin{proof}
For $t\ge1$, take a spine $v_0v_1\cdots v_t$.  Attach $a_1-1$ leaves to
$v_0$, attach $a_{i+1}-a_i-1$ leaves to $v_i$ for
$1\le i<t$, and attach $n-a_t-1$ leaves to $v_t$.  All multiplicities are
nonnegative.  Their sum is $n-t-1$, so the resulting tree has $n$ vertices.
Deleting the spine edge $v_{i-1}v_i$ leaves a prefix component of order
$a_i$.  Every other edge is a leaf edge.

There are $n-1-t$ leaf edges and $t$ spine edges.  Since
\[
  w_n(a)-w_n(1)
  =\C{n-1}{4}-\C{a}{4}-\C{n-a}{4}
  =\Delta_n(a),
\]
subtracting the star contribution and adding the $t$ spine corrections gives
\eqref{eq:caterpillar-formula}.
\end{proof}

Consequently, for fixed $n$, every subset sum of
$\Delta_n(2),\ldots,\Delta_n(n-2)$ is realized by a concrete caterpillar.
This statement is used repeatedly below; it is not merely a numerical encoding
of cut sizes.

\section{The exact all-tree certificate}

This section describes the finite computation that determines every exception.
Its completeness rests on a rooted-tree recurrence, not on a database of
selected tree shapes.

For subsets $X,Y\subseteq\N\cup\{0\}$, write
$X+Y=\{x+y:x\in X,\ y\in Y\}$.
Fix a global tree order $n$.  Let $P_n(r)$ be the set of internal edge sums of
rooted trees of order $r$, where the edge from the root to a possible parent is
omitted and every internal edge is evaluated with the global weight $w_n$.
Let $F_n(t)$ be the corresponding set for ordered forests of rooted trees with
total order $t$.

\begin{proposition}[Rooted-forest recurrence]\label{prop:rooted-dp}
For fixed $n$,
\begin{align}
  F_n(0)&=\{0\}, & P_n(1)&=\{0\},\label{eq:dp-base}\\
  P_n(r)&=F_n(r-1) &&(r\ge2),\label{eq:dp-root}\\
  F_n(t)&=\bigcup_{s=1}^{t}
  \left(F_n(t-s)+w_n(s)+P_n(s)\right) &&(t\ge1).
  \label{eq:dp-rec}
\end{align}
Moreover, $P_n(n)$ is exactly the set of values $\SW(T)$ over all
$n$-vertex trees $T$.
\end{proposition}

\begin{proof}
A rooted tree of order $r$ consists of its root together with a forest of
rooted child components of total order $r-1$, proving \eqref{eq:dp-root}.
For an ordered forest of total order $t$, choose its first rooted component to
have order $s$.  Its attachment edge contributes $w_n(s)$ and its internal
edges contribute an element of $P_n(s)$; the remaining ordered forest
contributes an element of $F_n(t-s)$.  This proves that every construction on
the right side of \eqref{eq:dp-rec} is valid.

Conversely, every rooted forest may be ordered arbitrarily, and after choosing
its first component it appears in one of the terms on the right.  Thus the
recurrence omits no rooted forest.  Ordering children may duplicate a value,
but creates no invalid value.  Rooting any unrooted $n$-vertex tree at an
arbitrary vertex does not change its edge contributions, so $P_n(n)$ is the
exact all-tree value set.
\end{proof}

\subsection{Bitset implementation and finite cutoffs}

For a cap $B$, a set $X\subseteq\{0,\ldots,B\}$ is encoded by the Python
integer
\[
  \beta(X)=\sum_{x\in X}2^x.
\]
A shift by $c$ represents $X+c$.  The Minkowski sum $X+Y$ is computed exactly
by iterating over the set bits of the smaller operand and OR-ing shifted copies
of the other operand.  Bits above $B$ are discarded.  Since every edge weight
is nonnegative, discarding a partial sum above $B$ cannot remove a later sum at
most $B$.

\begin{certificate}[Exceptional range]\label{cert:exceptions}
The program \code{sw4_certificate.py}, implementing
Proposition~\ref{prop:rooted-dp} with arbitrary-precision Python integers,
certifies the following statements.
\begin{enumerate}[label=(\alph*)]
\item Exactly $14{,}099$ positive integers at most $80{,}163$ do not belong to
      $\Tfour$.
\item The integer $80{,}163$ does not belong to $\Tfour$.
\item Every integer in $\interval{80{,}164}{178{,}857}$ belongs to $\Tfour$.
\end{enumerate}
\end{certificate}

\begin{proof}[Why the finite computation is exhaustive]
For the first two assertions, Corollary~\ref{cor:star-min} gives
\[
  G(29)=91{,}728>80{,}163.
\]
Hence no tree of order at least $29$ can contribute a value in the exceptional
range.  The program computes $P_n(n)\cap[0,80{,}163]$ exactly for every
$4\le n\le28$ and takes their union.

For the third assertion,
\[
  G(34)=180{,}048>178{,}857,
\]
so it is enough to compute the exact sets for $4\le n\le33$ up to the cap
$178{,}857$.  The program asserts that the missing list up to this larger cap
is identical to the $14{,}099$-element list from the first computation.  The
list is stored in \code{sw4_missing_values.txt}; the maximal consecutive
ranges are stored in \code{sw4_missing_ranges.txt}.
\end{proof}

The same source also performs an exact caterpillar subset-sum computation.
Since
\[
  G(81)=6{,}572{,}800,
\]
orders at most $80$ suffice below $G(81)$.  Applying
Lemma~\ref{lem:caterpillar-realization} to all subsets of
$\{2,\ldots,n-2\}$ for $4\le n\le80$ gives the next certificate.

\begin{certificate}[Initial caterpillar interval]\label{cert:initial-cat}
Every integer in
\[
  \interval{178{,}858}{6{,}572{,}799}
\]
belongs to $\Tfour$.
\end{certificate}

For fixed $n$, the subset sums are encoded by the exact update
\[
  B\longleftarrow B\ \mathrm{OR}\ (B\ll\Delta_n(a)),
  \qquad a=2,\ldots,n-2,
\]
with a mask at the stated cap.  This enumerates every subset and no other
choice.

\section{Central moment switches}

We next develop the common algebra behind the finite compressed bridge and the
universal tail.

\subsection{Even-order coordinates}

Put $n=2h$ and define, for $0\le x\le h-2$,
\begin{equation}\label{eq:delta-def}
  \delta_h(x)=\Delta_{2h}(h+x).
\end{equation}
The distance $x=0$ corresponds to the single cut size $h$.  Every positive
$x$ corresponds to two distinct cut sizes $h-x$ and $h+x$; because
$\Delta_{2h}(h-x)=\Delta_{2h}(h+x)$, these are two labeled slots of the same
weight $\delta_h(x)$.

Expanding \eqref{eq:caterpillar-formula} gives
\begin{equation}\label{eq:delta-quartic}
\begin{split}
  \delta_h(x)={}&-\frac{x^4}{12}
  +\left(-\frac{h^2}{2}+\frac{3h}{2}-\frac{11}{12}\right)x^2\\
  &+\frac{7h^4}{12}-\frac{17h^3}{6}
  +\frac{59h^2}{12}-\frac{11h}{3}+1.
\end{split}
\end{equation}
Thus $\delta_h$ is an even quartic in $x$.

For a finite multiset $S$ of nonnegative distances, write
\[
  \Phi_h(S)=\sum_{x\in S}\delta_h(x),
\]
with multiplicities.

\begin{lemma}[Moment switch]\label{lem:moment-switch}
Let $P,Q$ be finite multisets satisfying
\[
  |P|=|Q|,
  \qquad
  \sum_{x\in P}x^2=\sum_{x\in Q}x^2.
\]
Then
\begin{equation}\label{eq:moment-difference}
  \Phi_h(P)-\Phi_h(Q)
  =-\frac1{12}
  \left(\sum_{x\in P}x^4-\sum_{x\in Q}x^4\right),
\end{equation}
which is independent of $h$.
\end{lemma}

\begin{proof}
In \eqref{eq:delta-quartic}, equal cardinalities cancel the constant term and
equal square sums cancel the coefficient of $x^2$.  The displayed fourth-power
term is the only term that remains.
\end{proof}

A \emph{state group} is a finite collection of multisets satisfying the two
common-moment conditions in Lemma~\ref{lem:moment-switch}.  Its normalized
increments are its state values minus their minimum.

\begin{lemma}[Independent slot allocation]\label{lem:slot-allocation}
For a state group $\mathcal S$, let
\[
  r_{\mathcal S}(x)=\max_{S\in\mathcal S}m_S(x),
\]
where $m_S(x)$ is the multiplicity of $x$ in $S$.  Suppose a collection of
state groups satisfies
\[
  \sum_{\mathcal S}r_{\mathcal S}(0)\le1,
  \qquad
  \sum_{\mathcal S}r_{\mathcal S}(x)\le2\quad(x>0).
\]
Then all group choices can be made independently in a single caterpillar.
Every labeled slot not reserved by a group may also be made independently
optional.
\end{lemma}

\begin{proof}
Label the one slot at distance $0$ and the two slots at each positive distance.
For each group and distance, reserve $r_{\mathcal S}(x)$ labeled slots, with
reservations disjoint across groups.  Every state uses no more than its
reservation.  A chosen collection of states therefore selects a set of
distinct cut sizes.  Lemma~\ref{lem:caterpillar-realization} realizes their
union as a caterpillar.  Unreserved slots are disjoint from every state and may
be selected or omitted independently.
\end{proof}

\begin{lemma}[Bounded complete sequence]\label{lem:complete-sequence}
Suppose independent choices already realize every integer in $[0,W]$.
If an additional independent binary switch has increment $d\le W+1$, then the
enlarged choices realize every integer in $[0,W+d]$.  More generally, $A$
independent copies of that switch extend the interval to $[0,W+Ad]$.
\end{lemma}

\begin{proof}
One copy gives the union $[0,W]\cup[d,d+W]$.  The hypothesis makes these
intervals overlap or touch.  Repeating the same argument $A$ times proves the
second statement.
\end{proof}

\subsection{The universal and full seeds}

The first state group is
\begin{align*}
  S_0&=(0,2,2,3,6,6,10,10,12,12),\\
  S_1&=(1,1,4,4,6,7,8,9,12,13),\\
  S_2&=(1,1,4,5,5,7,7,11,11,13).
\end{align*}
All three have cardinality $10$ and square sum $577$; their fourth-power sums
are $64{,}177$, $64{,}165$, and $64{,}153$.  By
Lemma~\ref{lem:moment-switch}, their normalized increments are $0,1,2$.
The independent pair
\begin{align*}
  X&=(16,22,24,27,30,32,36),\\
  Y&=(17,20,25,26,31,33,35)
\end{align*}
has cardinality $7$, square sum $5{,}265$, and fourth-power sums
$4{,}701{,}201$ and $4{,}701{,}189$, so it has increment $1$.  These two
groups realize $[0,3]$.

Appendix~\ref{app:seed-witnesses} lists fifteen further binary pairs with
increments
\[
  4,8,16,\ldots,65{,}536.
\]
Their simultaneous demand is at most one at distance zero and at most two at
every positive distance; Appendix~\ref{app:capacity} records the compact
capacity audit.  Successively applying Lemma~\ref{lem:complete-sequence}
gives the universal seed interval
\begin{equation}\label{eq:universal-seed}
  [0,131{,}071].
\end{equation}
Its largest distance is $99$.

For finite junction calculations only, one more pair is used:
\[
  (32,52,76,100),
  \qquad
  (36,44,84,96).
\]
Both states have square sum $19{,}504$, and their fourth-power sums are
$141{,}722{,}368$ and $140{,}149{,}504$, giving increment $131{,}072$.
Together with \eqref{eq:universal-seed}, this produces the full seed
\begin{equation}\label{eq:full-seed}
  [0,262{,}143].
\end{equation}
Its largest distance is $100$, and its full simultaneous capacity is valid.

\section{Exact finite bridges}

\subsection{Bitset bridge through $292{,}031{,}545$}

For each $70\le n\le105$, the program
\code{sw4_caterpillar_bridge_certificate.py} computes every bounded subset
sum of $\Delta_n(2),\ldots,\Delta_n(n-2)$, shifts by $G(n)$, and ORs the
results over all orders.  The cap is $292{,}031{,}545$.

\begin{certificate}[Lower caterpillar bridge]\label{cert:lower-bridge}
The exact bit mask contains every bit in
\[
  \interval{6{,}572{,}800}{292{,}031{,}545}.
\]
Therefore this interval is contained in $\Tfour$.
\end{certificate}

The update is the standard exact subset recurrence from
Certificate~\ref{cert:initial-cat}.  There is no sampling, hashing, modular
filter, or floating-point arithmetic.

\subsection{Exact compressed interval unions}

The next certificate uses the fact that a finite subset of $\mathbb Z$ can be
stored exactly as a union of maximal integer intervals.

\begin{lemma}[Exact optional-weight update]\label{lem:interval-update}
Let $S\subseteq\mathbb Z$ be represented exactly by pairwise-disjoint maximal
integer intervals.  For an optional weight $w\ge0$, the attainable set after
allowing the choices $0$ and $w$ is
\[
  S'=S\cup(S+w).
\]
Merging the two sorted interval lists and coalescing only overlapping or
adjacent intervals produces the exact maximal-interval representation of
$S'$.
\end{lemma}

\begin{proof}
Translation sends each interval $[a,b]$ of $S$ to $[a+w,b+w]$ and no other
integer.  The union of the original and translated lists is therefore exactly
$S'$.  Coalescing overlapping or adjacent intervals changes only its
representation, not the represented set.
\end{proof}

For a fixed half-order $h$, the certificate constructs a baseline consisting
of $G(2h)$ plus the minimum value from every reserved state group.  The seed
state choices give an initial interval $[0,W_{\rm seed}]$.  Every annulus
switch and every remaining labeled cut slot supplies an independent optional
weight.  Lemma~\ref{lem:interval-update} is applied to every weight.  Thus the
final interval union is the exact attainable increment set for the stated
caterpillar family.

Three seed families are used.  Their complete states are listed in
Appendix~\ref{app:finite-seeds}; the certificate checks equal moments,
normalized increments, support validity, and simultaneous capacities before
processing any order.

\begin{table}[ht]
\centering
\caption{Families used in the compressed finite bridge.  The displayed first
and last intervals are the selected longest components at the endpoints of
each order range.}
\label{tab:finite-bridge-families}
\small
\begin{tabular}{@{}lcccc@{}}
\toprule
Family & $h$ range & Seed width & Fixed annuli & Endpoint intervals \\
\midrule
$\mathcal F_{51}$ & $63$--$66$ & $63$ & none &
$[282{,}204{,}658,495{,}014{,}629]$;\\
&&&&$[346{,}215{,}511,635{,}465{,}903]$\\[2pt]
$\mathcal F_{65}$ & $67$--$206$ & $2{,}047$ & none &
$[494{,}227{,}243,564{,}426{,}072]$;\\
&&&&$[58{,}222{,}371{,}640,236{,}605{,}611{,}082]$\\[2pt]
$\mathcal F_{100}$ & $206$--$253$ & $262{,}143$ & $(101,5),(141,8)$ &
$[106{,}740{,}323{,}728,189{,}108{,}214{,}894]$;\\
&&&&$[272{,}308{,}250{,}181,554{,}860{,}689{,}583]$\\
\bottomrule
\end{tabular}
\end{table}

The full output contains $192$ order rows.  Each selected interval begins no
later than one past the previously covered endpoint.  The first interval lies
inside Certificate~\ref{cert:lower-bridge}; every family transition also
overlaps.

\begin{certificate}[Compressed finite bridge]\label{cert:compressed-bridge}
The program \code{sw4_complete_bridge_certificate.cpp} and its exact output
\code{sw4_complete_bridge_output.txt} certify
\[
  \interval{282{,}204{,}658}{554{,}860{,}689{,}583}
  \subseteq\Tfour.
\]
Together with Certificate~\ref{cert:lower-bridge},
\begin{equation}\label{eq:finite-bridge}
  \boxed{
  \interval{6{,}572{,}800}{554{,}860{,}689{,}583}
  \subseteq\Tfour.}
\end{equation}
\end{certificate}

The independent Python checker
\code{sw4_bridge_independent_check.py} reads all $192$ rows, verifies every
junction, and recomputes the six boundary orders independently.

\begin{remark}[Necessary correction to a weaker propagation argument]
It is essential that the certificate use the exact interval update of
Lemma~\ref{lem:interval-update}.  At $h=206$, after a partial insertion, a raw
weight $12{,}794{,}200$ exceeds the then-current complete width plus one,
$12{,}772{,}104$.  Hence the sufficient condition of
Lemma~\ref{lem:complete-sequence} does not propagate a single interval through
that ordering.  The exact attainable set temporarily has more than one
component; subsequent weights reconnect the components.  The compressed
interval algorithm retains every component and therefore proves the stated
interval without the invalid stronger claim.
\end{remark}

\section{Thue--Morse annuli}

We now prove the universal infinite tail.  The construction uses the universal
seed \eqref{eq:universal-seed}, not the extra $131{,}072$ digit.

\subsection{An eight-point identity}

For integers $R,q\ge1$ and $0\le r<q$, set
\[
  x_j=R+r+jq,
  \qquad 0\le j\le7.
\]
Partition the indices according to the parity of the number of ones in the
binary expansion of $j$.  Equivalently, put $\varepsilon_j=1$ for
$j\in\{0,3,5,6\}$ and $\varepsilon_j=-1$ for
$j\in\{1,2,4,7\}$.

A direct summation gives
\begin{equation}\label{eq:tm-moments}
\begin{array}{c|rrrrr}
 m&0&1&2&3&4\\ \hline
 \sum_{j=0}^{7}\varepsilon_jj^m&0&0&0&-48&-672.
\end{array}
\end{equation}
Let $E_{R,q,r}$ and $O_{R,q,r}$ denote the two four-element distance states.

\begin{lemma}[Thue--Morse switch]\label{lem:tm-switch}
The two states have an index difference independent of $h$, of absolute value
\begin{equation}\label{eq:tm-digit}
  d_{R,q,r}=8q^3\bigl(2(R+r)+7q\bigr).
\end{equation}
For fixed $R,q$, the $q$ switches $r=0,\ldots,q-1$ have disjoint supports,
and their supports partition the annulus of distances
\[
  R,R+1,\ldots,R+8q-1.
\]
Their total increment is
\begin{equation}\label{eq:annulus-total}
  D(R,q)=8q^4(2R+8q-1).
\end{equation}
\end{lemma}

\begin{proof}
Write $t=R+r$.  In the signed difference of the two state values, the constant
and quadratic parts of \eqref{eq:delta-quartic} vanish by
\eqref{eq:tm-moments}.  Expanding the fourth power gives
\[
  \sum_{j=0}^{7}\varepsilon_j(t+jq)^4
  =4tq^3(-48)+q^4(-672).
\]
Multiplication by the coefficient $-1/12$ in
\eqref{eq:delta-quartic} yields
\[
  16tq^3+56q^4=8q^3(2t+7q),
\]
which is \eqref{eq:tm-digit}.

Every integer in $[R,R+8q-1]$ has a unique expression $R+r+jq$ with
$0\le r<q$ and $0\le j\le7$, proving the support assertion.  Finally,
\begin{align*}
  \sum_{r=0}^{q-1}d_{R,q,r}
  &=8q^3\sum_{r=0}^{q-1}\bigl(2R+2r+7q\bigr)\\
  &=8q^4(2R+8q-1).
\end{align*}
\end{proof}

\section{A universal fixed-order interval}

\subsection{The annulus ladder}

Define
\begin{equation}\label{eq:ladder}
  R_0=100,
  \qquad q_0=4,
  \qquad
  R_{i+1}=R_i+8q_i,
  \qquad
  q_{i+1}=\left\lceil\frac{3q_i}{2}\right\rceil.
\end{equation}
Let
\begin{equation}\label{eq:Wk}
  W_0=131{,}071,
  \qquad
  W_{k+1}=W_k+D(R_k,q_k).
\end{equation}
Thus $W_k$ is the complete switch width after the universal seed and the first
$k$ annuli, indexed $0,\ldots,k-1$.

\begin{lemma}[Annulus ladder completeness]\label{lem:annulus-complete}
For every $k\ge0$, the universal seed together with the first $k$ annuli
realizes every increment in $[0,W_k]$.
\end{lemma}

\begin{proof}
For the first three annuli, the first digits and preceding widths are
\[
  116{,}736\le131{,}072,
  \qquad
  528{,}768\le604{,}160,
  \qquad
  2{,}466{,}936\le3{,}828{,}608.
\]
Thus Lemma~\ref{lem:complete-sequence} starts the induction.

We first prove the invariant
\begin{equation}\label{eq:Rq-invariant}
  12q_i\le R_i\le25q_i.
\end{equation}
It holds for $(R_0,q_0)=(100,4)$.  Since $q_i\ge4$,
\[
  \frac32q_i\le q_{i+1}\le\frac{13}{8}q_i.
\]
Consequently,
\[
  R_{i+1}\le33q_i\le22q_{i+1}<25q_{i+1},
\]
and
\[
  R_{i+1}\ge20q_i\ge\frac{160}{13}q_{i+1}>12q_{i+1}.
\]

Consider an annulus with index $i\ge3$.  Put $p=q_{i-1}$ and
$R=R_{i-1}$.  Then $p\ge9$, $R/p\ge12$, and
$q_i/p\le14/9$.  Since $R_i=R+8p$, equation
\eqref{eq:tm-digit} gives
\begin{equation}\label{eq:first-digit-upper}
  \frac{d_{R_i,q_i,0}}{8p^4}
  \le
  \left(\frac{14}{9}\right)^3
  \left(2\frac Rp+\frac{242}{9}\right).
\end{equation}
The preceding annulus has total increment
\begin{equation}\label{eq:previous-annulus-lower}
  \frac{D(R,p)}{8p^4}
  =p\left(2\frac Rp+8-\frac1p\right).
\end{equation}
The difference between the right side of
\eqref{eq:previous-annulus-lower} and the right side of
\eqref{eq:first-digit-upper} is increasing in $R/p$.  Its minimum in the
stated domain occurs at $p=9$ and $R/p=12$, where the required inequality is
the exact integer comparison
\[
  287\cdot6561>2744\cdot458.
\]
Therefore
\[
  d_{R_i,q_i,0}\le D(R_{i-1},q_{i-1})\le W_i.
\]

Within a fixed annulus,
\[
  d_{R,q,r+1}-d_{R,q,r}=16q^3,
\]
and $d_{R,q,r+1}\le2d_{R,q,r}$.  If the current width before inserting
$d_{R,q,r}$ is at least $d_{R,q,r}-1$, then afterward it is at least
$2d_{R,q,r}-1$, so the next digit also satisfies the hypothesis of
Lemma~\ref{lem:complete-sequence}.  This proves the induction over all digits
and all annuli.
\end{proof}

\subsection{Reaching the raw caterpillar cuts}

For $h\ge365$, let $k=k(h)$ be the largest integer satisfying
\begin{equation}\label{eq:k-choice}
  R_k\le h-1.
\end{equation}
Then the first $k$ annuli occupy distances $100$ through $R_k-1$ and fit
inside $0\le x\le h-2$.  Maximality gives $h\le R_{k+1}$.
The smallest nonzero raw caterpillar increment is
\begin{equation}\label{eq:delta2}
  \Delta_{2h}(2)=\C{2h-2}{3}.
\end{equation}

\begin{lemma}[Switch width reaches the first raw cut]\label{lem:width-raw}
For every $h\ge365$,
\begin{equation}\label{eq:W-raw}
  W_{k(h)}\ge\Delta_{2h}(2)-1.
\end{equation}
\end{lemma}

\begin{proof}
The first relevant ladder values are
\begin{center}
\begin{tabular}{@{}r r r@{}}
\toprule
$k$ & $W_k$ & $\binom{2R_{k+1}-2}{3}-1$\\
\midrule
$4$ & $215{,}457{,}655$ & $199{,}064{,}819$\\
$5$ & $1{,}607{,}941{,}615$ & $648{,}686{,}323$\\
$6$ & $12{,}672{,}515{,}567$ & $2{,}138{,}222{,}579$\\
\bottomrule
\end{tabular}
\end{center}
These are exact evaluations of \eqref{eq:ladder}--\eqref{eq:Wk}.

For $k\ge7$, put $i=k-1$ and $p=q_i$.  Then $p\ge48$.  From
\eqref{eq:Rq-invariant} and $q_{i+1}\le13p/8$,
\[
  R_{i+2}=R_i+8p+8q_{i+1}\le46p.
\]
Moreover,
\[
  W_k\ge D(R_i,p)
  =8p^4(2R_i+8p-1)
  \ge64p^5.
\]
On the other hand,
\[
  \C{2R_{i+2}-2}{3}
  <\frac{(2R_{i+2})^3}{6}
  \le\frac43\,46^3p^3.
\]
The inequality
\[
  64\cdot48^2>\frac43\,46^3
\]
therefore implies
$W_k\ge\binom{2R_{k+1}-2}{3}-1$.  Since
$h\le R_{k+1}$, monotonicity of the binomial coefficient and
\eqref{eq:delta2} prove \eqref{eq:W-raw}.
\end{proof}

\subsection{Insertion of all remaining slots}

For $1\le a<n/2$, equation \eqref{eq:w-difference} implies that the first
differences of $\Delta_n(a)$ are positive and nonincreasing.  Since
$\Delta_n(1)=0$, we obtain
\begin{equation}\label{eq:concavity-bound}
  \Delta_n(a)\le(a-1)\Delta_n(2),
  \qquad 2\le a\le n/2.
\end{equation}

The seed uses only distances at most $99$.  Each annulus uses one of the two
slots at each supported positive distance.  Therefore, for every
$a=2,\ldots,h-100$, at least one slot of weight $\Delta_{2h}(a)$ remains
unreserved: it is the slot at distance $x=h-a\ge100$.

Starting from the complete interval $[0,W_k]$, insert one such slot in the
order $a=2,3,\ldots,h-100$.  Lemma~\ref{lem:width-raw} handles $a=2$.
After the slots through $a$ have been inserted, the current width $W$ satisfies
\begin{align*}
  W+1
  &\ge \Delta_{2h}(2)+\sum_{j=2}^{a}\Delta_{2h}(j)\\
  &\ge a\Delta_{2h}(2)\\
  &\ge\Delta_{2h}(a+1),
\end{align*}
where the last line uses \eqref{eq:concavity-bound}.  Thus the interval remains
complete.

Let
\begin{equation}\label{eq:Sh-def}
  S_h=\sum_{a=2}^{h-100}\Delta_{2h}(a).
\end{equation}
Hockey-stick summation gives
\begin{equation}\label{eq:Sh-closed}
\begin{split}
  S_h={}&(h-101)\C{2h-1}{4}-\C{h-99}{5}\\
       &-\C{2h-1}{5}+\C{h+100}{5}.
\end{split}
\end{equation}
Furthermore, exact expansion yields
\begin{equation}\label{eq:Sh-central-poly}
\begin{split}
  120\bigl(S_h-\Delta_{2h}(h)\bigr)
  ={}&48h^5-7315h^4+34790h^3\\
    &+19641055h^2-59058458h+19539439680.
\end{split}
\end{equation}
For $h\ge365$, the first two terms combine to
$h^4(48h-7315)>0$, the terms
$19641055h^2-59058458h$ combine to a positive quantity, and every other
term displayed in \eqref{eq:Sh-central-poly} is positive.  Hence
\begin{equation}\label{eq:S-dominates-center}
  S_h\ge\Delta_{2h}(h).
\end{equation}

After the first copy at every outer size has been inserted, the accumulated
width is at least $S_h$, hence at least the largest raw cut weight by
\eqref{eq:S-dominates-center}.  Every remaining unreserved slot may therefore
be inserted in nondecreasing order using
Lemma~\ref{lem:complete-sequence}.

Define $L_h$ explicitly as
\begin{equation}\label{eq:Lh-def}
\begin{split}
 L_h={}&G(2h)
 +\sum_{\mathcal S\in\mathcal G}
   \min_{S\in\mathcal S}\Phi_h(S)\\
 &+\sum_{i=0}^{k(h)-1}\sum_{r=0}^{q_i-1}
   \min\bigl\{\Phi_h(E_{R_i,q_i,r}),
               \Phi_h(O_{R_i,q_i,r})\bigr\},
\end{split}
\end{equation}
where $\mathcal G$ is the universal seed.  Let $U_h$ be $L_h$ plus the
complete switch width and the weights of all remaining optional slots.

\begin{theorem}[Universal fixed-order interval]\label{thm:fixed-order}
For every integer $h\ge365$, the caterpillars described above realize every
integer in
\[
  J_h=[L_h,U_h].
\]
At the first order,
\begin{equation}\label{eq:J365}
  J_{365}=
  [1{,}596{,}122{,}112{,}921,
   3{,}577{,}546{,}733{,}067].
\end{equation}
\end{theorem}

\begin{proof}
The seed and annulus states are simultaneously valid by
Lemma~\ref{lem:slot-allocation}; their choices realize $[0,W_{k(h)}]$ by
Lemma~\ref{lem:annulus-complete}.  The preceding insertion argument adds every
remaining labeled slot while preserving a complete interval.  Translating by
the baseline \eqref{eq:Lh-def} proves the theorem.  The endpoint evaluation at
$h=365$ is an exact integer calculation verified independently by the supplied
Python certificates.
\end{proof}

\section{Overlap of the universal intervals}

\subsection{Finite overlap range}

The exact program \code{sw4_infinite_tail_certificate.py} evaluates the
explicit construction of Theorem~\ref{thm:fixed-order} for every
$365\le h\le1171$.  It verifies
\begin{equation}\label{eq:finite-universal-overlap}
  L_{h+1}\le U_h+1,
  \qquad 365\le h\le1170.
\end{equation}
The minimum overlap margin $U_h+1-L_{h+1}$ is
$1{,}960{,}269{,}091{,}932$, at the transition $h=365$ to $366$.
The $807$ exact rows are stored in
\code{sw4_infinite_tail_overlaps.csv}.

\subsection{Symbolic overlap range}

At order $2(h+1)$, every selected universal seed state contains exactly
$77$ cuts: ten from the three-state group, seven from the unit pair, and four
from each of the fifteen binary groups.  If $k=k(h+1)$ annuli are selected,
they use $4\sum_{i<k}q_i$ cuts.  From
$R_k=100+8\sum_{i<k}q_i$ and $R_k\le h$, the total number of selected
baseline cuts is at most
\begin{equation}\label{eq:cut-count-bound}
  77+\frac{h-100}{2}=\frac{h+54}{2}.
\end{equation}
Every selected cut at order $2h+2$ contributes at most
$\binom{2h+1}{4}$.  Therefore
\begin{equation}\label{eq:L-upper}
  L_{h+1}
  \le G(2h+2)+\frac{h+54}{2}\C{2h+1}{4}.
\end{equation}
On the other hand, Theorem~\ref{thm:fixed-order} and the insertion of the
outer cuts give
\begin{equation}\label{eq:U-lower}
  U_h\ge G(2h)+S_h.
\end{equation}
Substituting \eqref{eq:Sh-closed} into
\eqref{eq:L-upper}--\eqref{eq:U-lower} and clearing denominators gives
\begin{equation}\label{eq:overlap-poly}
\begin{split}
120\Biggl(&G(2h)+S_h+1-G(2h+2)
-\frac{h+54}{2}\C{2h+1}{4}\Biggr)\\
={}&8h^5-9365h^4+35340h^3+19643615h^2\\
&-59060078h+19539440040.
\end{split}
\end{equation}
For $h\ge1171$, the first two terms combine to
$h^4(8h-9365)>0$, the pair
$19643615h^2-59060078h$ is positive, and the remaining displayed terms are
positive.  Hence \eqref{eq:overlap-poly} is positive and
\begin{equation}\label{eq:symbolic-overlap}
  L_{h+1}\le U_h+1,
  \qquad h\ge1171.
\end{equation}

Combining \eqref{eq:finite-universal-overlap} and
\eqref{eq:symbolic-overlap} proves the universal overlap theorem.

\begin{theorem}[Universal tail above $J_{365}$]\label{thm:universal-tail}
\[
  \bigcup_{h\ge365}J_h
  =\interval{1{,}596{,}122{,}112{,}921}{\infty}.
\]
\end{theorem}

\begin{proof}
The first interval is \eqref{eq:J365}.  Every consecutive pair touches or
overlaps by \eqref{eq:finite-universal-overlap} and
\eqref{eq:symbolic-overlap}.  The left endpoints tend to infinity because
$L_h\ge G(2h)$, so the union is exactly the displayed ray.
\end{proof}

\section{Exact finite junction to the universal tail}

The first universal interval begins above the finite bridge
\eqref{eq:finite-bridge}.  A final exact chain uses the full seed
\eqref{eq:full-seed}, the fixed annuli
\begin{equation}\label{eq:fixed-tail-annuli}
  (R,q)=(101,5),
  \qquad
  (R,q)=(141,8),
\end{equation}
and every remaining labeled cut slot.
For each $h=253,254,\ldots,365$, the C++ program
\code{sw4_infinite_tail_finite_overlap.cpp} applies the exact interval
update of Lemma~\ref{lem:interval-update} and selects a longest component.

The first selected interval, at $h=253$, is
\begin{equation}\label{eq:tail-junction-first}
  [272{,}308{,}250{,}181,
   554{,}860{,}689{,}583],
\end{equation}
which already lies in and reaches the end of the finite bridge.  Every one of
the $113$ selected intervals overlaps the preceding covered set.  The last,
at $h=365$, is
\begin{equation}\label{eq:tail-junction-last}
  [1{,}304{,}696{,}148{,}031,
   3{,}868{,}972{,}697{,}957],
\end{equation}
which overlaps $J_{365}$.

\begin{certificate}[Finite junction]\label{cert:tail-junction}
The exact interval rows in
\code{sw4_infinite_tail_finite_overlap.csv} certify
\[
  \interval{272{,}308{,}250{,}181}
           {3{,}868{,}972{,}697{,}957}
  \subseteq\Tfour,
\]
and the final endpoint overlaps the first universal interval.
\end{certificate}

\begin{corollary}[Self-contained infinite tail]\label{cor:infinite-tail}
\begin{equation}\label{eq:tail-final}
  \boxed{
  \interval{272{,}308{,}250{,}181}{\infty}
  \subseteq\Tfour.}
\end{equation}
\end{corollary}

\begin{proof}
Certificate~\ref{cert:tail-junction} joins continuously to
Theorem~\ref{thm:universal-tail}.
\end{proof}

\section{Deduction of the exact inverse parameters}

We now combine the logically independent pieces.

\begin{proof}[Proof of Theorem~\ref{thm:main}]
Certificate~\ref{cert:exceptions} proves that exactly $14{,}099$ positive
integers at most $80{,}163$ are unattainable, that $80{,}163$ itself is
unattainable, and that every integer from $80{,}164$ through $178{,}857$ is
attainable.  Certificate~\ref{cert:initial-cat} continues the coverage through
$6{,}572{,}799$.  Equation \eqref{eq:finite-bridge} continues it through
$554{,}860{,}689{,}583$.  Finally, the infinite tail
\eqref{eq:tail-final} begins below that endpoint.  Hence every integer greater
than $80{,}163$ is attainable, while the certified $14{,}099$ exceptions at or
below it are the complete complement.  Therefore
\[
  e(4)=14{,}099,
  \qquad
  N_0(4)=80{,}163.
\]
\end{proof}

\section{Computational model, reproducibility, and audit}

\subsection{Arithmetic and algorithmic completeness}

All Python computations use arbitrary-precision integers.  The bitset
certificates therefore have no fixed-width overflow.  The C++ programs use
signed $128$-bit intermediates for binomial products and assert before every
conversion that stored endpoints fit in signed $64$-bit range.  No
floating-point arithmetic is used.

The roles of the programs are as follows.
\begin{description}[style=nextline]
\item[\code{sw4_certificate.py}]
Implements the exact rooted-forest recurrence and the original caterpillar
subset-sum certificate.  Its order bounds are justified by the star minimum,
so it covers every tree order that can affect the claimed ranges.
\item[\code{sw4_caterpillar_bridge_certificate.py}]
Computes all bounded caterpillar subset sums for orders $70$ through $105$ and
checks the entire required bit mask.
\item[\code{sw4_complete_bridge_certificate.cpp}]
Computes exact interval unions for the $192$ finite-bridge rows and asserts
every junction.
\item[\code{sw4_bridge_independent_check.py}]
Independently checks all bridge rows and recomputes the six family-boundary
orders.
\item[\code{sw4_infinite_tail_finite_overlap.cpp}]
Computes the $113$ exact finite-junction rows from $h=253$ through $365$.
\item[\code{sw4_infinite_tail_certificate.py}]
Checks every seed witness and capacity, all complete-sequence inequalities in
the universal construction, the explicit interval at $h=365$, and every
finite overlap through $h=1171$.
\item[\code{sw4_infinite_tail_independent_check.py}]
Independently checks the witness tables, both finite chains, all $807$
universal overlap rows, and the symbolic boundary identities.
\end{description}

The exact optional-weight update of Lemma~\ref{lem:interval-update} is exhaustive
because each step replaces the current exact set $S$ by precisely
$S\cup(S+w)$.  Induction on the number of optional weights proves that the
final union contains every and only the sums generated by the stated family.
Likewise, the bitset subset update is the exact recurrence for choosing or
omitting each cut.

\subsection{Reproduction commands}

From the supplementary archive directory, the principal checks are reproduced
by
\begin{lstlisting}[basicstyle=\ttfamily\small,frame=single]
python sw4_certificate.py
python sw4_caterpillar_bridge_certificate.py

g++ -O3 -std=c++20 -Wall -Wextra -pedantic \
  sw4_complete_bridge_certificate.cpp \
  -o sw4_complete_bridge_certificate
./sw4_complete_bridge_certificate
python sw4_bridge_independent_check.py

g++ -O3 -std=c++20 -Wall -Wextra -pedantic \
  sw4_infinite_tail_finite_overlap.cpp \
  -o sw4_infinite_tail_finite_overlap
./sw4_infinite_tail_finite_overlap
python sw4_infinite_tail_certificate.py
python sw4_infinite_tail_independent_check.py

sha256sum -c sw4_self_contained_sha256.txt
\end{lstlisting}
The expected summaries include
\begin{lstlisting}[basicstyle=\ttfamily\small,frame=single]
PASS
candidate e(4) = 14099
candidate N0(4) = 80163
exact all-tree coverage: 80164..178857
caterpillar coverage: 178858..6572799

PASS
orders=70..105
coverage=6572800..292031545

PASS
bridge_start=282204658
bridge_end=554860689583

PASS
bridge_start=272308250181
bridge_end=3868972697957
universal_L_365=1596122112921

PASS
universal_h_start=365
universal_first_interval=1596122112921..3577546733067
symbolic_overlap_h>=1171
\end{lstlisting}

The verified supplementary archive has SHA-256
\begin{center}
\code{1ed15206d460cac319b0b3683a41ade37db687ecd590a199126b08e189b47342}.
\end{center}
Selected file hashes appear in Appendix~\ref{app:hashes}; the full manifest is
\code{sw4_self_contained_sha256.txt}.  The supplementary archive is cited
as reference~\cite{supplement}.

\subsection{Hostile audit summary}

The following failure modes are explicitly excluded.
\begin{enumerate}[label=(\arabic*)]
\item Proposition~\ref{prop:rooted-dp} covers all rooted trees; the star lower
      bound excludes all uncomputed orders below the exceptional cutoff.
\item Every caterpillar state is a set of labeled cut slots satisfying the
      one-at-zero and two-at-positive-distance capacities; Lemma
      \ref{lem:caterpillar-realization} converts every selected set into an
      actual finite tree.
\item Every seed row has equal cardinality and equal square sum, and its
      fourth-power difference gives the claimed exact increment.
\item Thue--Morse supports are pairwise disjoint within an annulus and across
      the ladder; their increment formula is derived symbolically.
\item The universal complete-sequence inequalities are proved for all annuli,
      not inferred from numerical density.
\item The finite interval programs preserve every component.  In particular,
      the invalid stronger one-interval propagation noted after
      Certificate~\ref{cert:compressed-bridge} is not used.
\item Every finite interval junction is checked with exact integers.  The
      transition to the universal range is exact, and the symbolic overlap
      polynomial is positive throughout its entire claimed domain.
\item The final finite bridge and the infinite tail overlap, leaving no
      unproved integer above $80{,}163$.
\end{enumerate}

\appendix

\section{Universal central-seed witnesses}\label{app:seed-witnesses}

For each row below, the two states have the same cardinality and the same
square sum.  The normalized increment is the absolute fourth-power difference
divided by $12$, as in Lemma~\ref{lem:moment-switch}.  The orientation of the
states is irrelevant because the lower-valued state is used in the baseline.

\begin{landscape}
\scriptsize
\begin{longtable}{@{}r p{0.245\linewidth} p{0.245\linewidth} r r r@{}}
\caption{Binary witnesses for the universal central seed.}\label{tab:universal-witnesses}\\
\toprule
$d$ & State $P$ & State $Q$ & Common $\sum x^2$ & $\sum_{P}x^4$ & $\sum_{Q}x^4$\\
\midrule
\endfirsthead
\toprule
$d$ & State $P$ & State $Q$ & Common $\sum x^2$ & $\sum_{P}x^4$ & $\sum_{Q}x^4$\\
\midrule
\endhead
$4$ & $72,78,85,90$ & $73,76,88,88$ & $26{,}593$ & $181{,}699{,}537$ & $181{,}699{,}489$\\
$8$ & $8,58,71,73$ & $45,46,51,84$ & $13{,}798$ & $65{,}130{,}514$ & $65{,}130{,}418$\\
$16$ & $77,83,90,95$ & $78,81,93,93$ & $29{,}943$ & $229{,}671{,}987$ & $229{,}672{,}179$\\
$32$ & $81,87,94,99$ & $82,85,97,97$ & $32{,}767$ & $274{,}470{,}979$ & $274{,}471{,}363$\\
$64$ & $22,43,64,71$ & $25,48,54,75$ & $11{,}470$ & $45{,}841{,}954$ & $45{,}842{,}722$\\
$128$ & $15,28,39,51$ & $23,23,37,52$ & $5{,}131$ & $9{,}743{,}923$ & $9{,}745{,}459$\\
$256$ & $9,24,64,87$ & $29,40,40,91$ & $12{,}322$ & $74{,}405{,}314$ & $74{,}402{,}242$\\
$512$ & $14,14,29,35$ & $18,18,21,37$ & $2{,}458$ & $2{,}284{,}738$ & $2{,}278{,}594$\\
$1{,}024$ & $28,31,86,92$ & $30,33,80,96$ & $17{,}605$ & $127{,}878{,}289$ & $127{,}890{,}577$\\
$2{,}048$ & $16,60,66,66$ & $16,62,62,68$ & $12{,}568$ & $50{,}975{,}008$ & $50{,}999{,}584$\\
$4{,}096$ & $55,61,67,70$ & $57,58,69,69$ & $16{,}135$ & $67{,}157{,}587$ & $67{,}206{,}739$\\
$8{,}192$ & $26,41,53,65$ & $27,38,57,63$ & $9{,}391$ & $29{,}023{,}843$ & $28{,}925{,}539$\\
$16{,}384$ & $3,60,67,75$ & $19,53,68,77$ & $13{,}723$ & $64{,}751{,}827$ & $64{,}555{,}219$\\
$32{,}768$ & $13,20,43,47$ & $15,21,44,45$ & $4{,}627$ & $8{,}487{,}043$ & $8{,}093{,}827$\\
$65{,}536$ & $38,50,56,70$ & $42,46,54,72$ & $11{,}980$ & $42{,}179{,}632$ & $42{,}966{,}064$\\
\bottomrule
\end{longtable}
\end{landscape}

\section{Capacity audit for the central seeds}\label{app:capacity}

For the universal seed, the total reserved demand equals one at
\[
\begin{split}
\{&0,17,19,32,36,39,41,42,44,47,48,50,52,55,56,61,63,65,\\
  &76,80,82,83,84,86,91,92,94,95,96,99\},
\end{split}
\]
equals two at
\[
\begin{split}
\{&1,2,3,4,5,6,7,8,9,10,11,12,13,14,15,16,18,20,21,22,23,24,25,26,\\
  &27,28,29,30,31,33,35,37,38,40,43,45,46,51,53,54,57,58,60,62,64,\\
  &66,67,68,69,70,71,72,73,75,77,78,81,85,87,88,90,93,97\},
\end{split}
\]
and is zero at the other distances through $99$.  Thus the one-slot capacity
at zero and the two-slot capacity at every positive distance are respected.
The full seed adds the $131{,}072$ pair.  Its exact demand remains at most one
at zero and at most two elsewhere; its largest used distance is $100$.  The
machine-readable versions are \code{sw4_central_seed_capacity.csv} and
\code{sw4_full_seed_capacity.csv}.

\section{Compact seeds used in the finite bridge}\label{app:finite-seeds}

All families include the common three-state group $S_0,S_1,S_2$ and the unit
pair $X,Y$.  The following tables list the remaining binary pairs.  In each
row the common square sum and the fourth-power difference verify the displayed
increment through Lemma~\ref{lem:moment-switch}.

\begin{table}[ht]
\centering
\caption{Additional pairs for $\mathcal F_{51}$, yielding seed width $63$.}
\small
\begin{tabular}{@{}rllrrr@{}}
\toprule
$d$ & $P$ & $Q$ & $\sum x^2$ & $\sum_Px^4$ & $\sum_Qx^4$\\
\midrule
$4$ & $23,28,32,49$ & $18,25,42,45$ & $4{,}738$ & $7{,}707{,}874$ & $7{,}707{,}922$\\
$8$ & $16,24,45,51$ & $8,28,47,49$ & $5{,}458$ & $11{,}263{,}138$ & $11{,}263{,}234$\\
$16$ & $20,29,46,50$ & $14,36,42,51$ & $5{,}857$ & $11{,}594{,}737$ & $11{,}594{,}929$\\
$32$ & $9,39,39,46$ & $23,29,37,50$ & $5{,}239$ & $9{,}110{,}899$ & $9{,}111{,}283$\\
\bottomrule
\end{tabular}
\end{table}

\begin{table}[ht]
\centering
\caption{Additional pairs for $\mathcal F_{65}$, yielding seed width $2{,}047$.}
\small
\begin{tabular}{@{}rllrrr@{}}
\toprule
$d$ & $P$ & $Q$ & $\sum x^2$ & $\sum_Px^4$ & $\sum_Qx^4$\\
\midrule
$4$ & $23,28,32,49$ & $18,25,42,45$ & $4{,}738$ & $7{,}707{,}874$ & $7{,}707{,}922$\\
$8$ & $21,24,45,51$ & $19,27,43,52$ & $5{,}643$ & $11{,}392{,}083$ & $11{,}392{,}179$\\
$16$ & $14,38,55,55$ & $3,43,54,54$ & $7{,}690$ & $20{,}424{,}802$ & $20{,}424{,}994$\\
$32$ & $9,39,39,46$ & $23,29,37,50$ & $5{,}239$ & $9{,}110{,}899$ & $9{,}111{,}283$\\
$64$ & $17,28,44,59$ & $15,36,37,60$ & $6{,}490$ & $16{,}563{,}634$ & $16{,}564{,}402$\\
$128$ & $18,35,48,48$ & $26,29,44,52$ & $6{,}157$ & $12{,}222{,}433$ & $12{,}223{,}969$\\
$256$ & $20,40,58,65$ & $16,42,60,63$ & $9{,}589$ & $31{,}887{,}121$ & $31{,}890{,}193$\\
$512$ & $57,57,63,64$ & $56,59,61,65$ & $14{,}563$ & $53{,}642{,}179$ & $53{,}648{,}323$\\
$1{,}024$ & $14,34,41,53$ & $15,30,46,51$ & $5{,}842$ & $12{,}090{,}994$ & $12{,}103{,}282$\\
\bottomrule
\end{tabular}
\end{table}

For $\mathcal F_{51}$ the total reservation uses $33$ selected cuts, has
maximum distance $51$, and respects all capacities.  For $\mathcal F_{65}$ it
uses $53$ selected cuts, has maximum distance $65$, and respects all
capacities.  The corresponding assertions are executed before the bridge
program processes any interval row.  The family $\mathcal F_{100}$ uses the
full universal table in Appendix~\ref{app:seed-witnesses} together with the
additional $131{,}072$ pair.

\section{Selected exact outputs and hashes}\label{app:hashes}

The complete interval tables contain $192$, $113$, and $807$ rows,
respectively.  They are supplied in full rather than typeset here.  The key
junctions are
\begin{center}
\small
\begin{tabularx}{\textwidth}{@{}l>{\raggedright\arraybackslash}X>{\raggedright\arraybackslash}X@{}}
\toprule
Certificate & First relevant interval & Final relevant interval\\
\midrule
Lower compressed bridge &
$[282{,}204{,}658,495{,}014{,}629]$ &
$[272{,}308{,}250{,}181,554{,}860{,}689{,}583]$\\
Finite tail junction &
$[272{,}308{,}250{,}181,554{,}860{,}689{,}583]$ &
$[1{,}304{,}696{,}148{,}031,3{,}868{,}972{,}697{,}957]$\\
Universal family &
$[1{,}596{,}122{,}112{,}921,3{,}577{,}546{,}733{,}067]$ &
all subsequent overlapping intervals\\
\bottomrule
\end{tabularx}
\end{center}

\begin{longtable}{@{}p{0.45\linewidth}p{0.49\linewidth}@{}}
\caption{Selected SHA-256 values from the verified manifest.}\\
\toprule
File & SHA-256\\
\midrule
\endfirsthead
\toprule
File & SHA-256\\
\midrule
\endhead
\code{sw4_certificate.py} &
\code{d85678eb9a4ecb97ea3e0786eb721bd6bb7dc407fbc012775eee76a469675e1d}\\
\code{sw4_missing_values.txt} &
\code{c5abe94bd3405d695e42b9e8970f9a95a6ac00e138f0959290cef19ca62c349f}\\
\code{sw4_missing_ranges.txt} &
\code{e70a93150c71066fcc6d64e585643b3acf989d0d335a5d90c79486cd3b9f115e}\\
\code{sw4_caterpillar_bridge_certificate.py} &
\code{02450eadd42abd1f182a2e2c562f2886e4703a6480bed76a4ef733408230946d}\\
\code{sw4_complete_bridge_certificate.cpp} &
\code{05ac482d6cafd8ec999ef2c087f706aae565cfd2b5292e6271724ce111d61aaf}\\
\code{sw4_complete_bridge_output.txt} &
\code{549cb209c95154e49c9c50876d3cec58bb332f994321cf44edb390a287664cd3}\\
\code{sw4_bridge_independent_check.py} &
\code{883242643f681627e167d78c9c3b54548b786cc75d8b227424f9d20230adcfea}\\
\code{sw4_infinite_tail_finite_overlap.cpp} &
\code{9c25622e5253f5ea3396bd66fd88b7885b2f414c1dc06ba3974360523b9e55e9}\\
\code{sw4_infinite_tail_finite_overlap.csv} &
\code{c3121d3dc41ebd5bcf16f7fd7fc452abdaac27d025c09d87e5542899f6d9258b}\\
\code{sw4_infinite_tail_certificate.py} &
\code{ebe6f805621a58f0284f0de54bbf36b433efe7a0a633c9283a6e9fb34445ee94}\\
\code{sw4_infinite_tail_overlaps.csv} &
\code{9c481a07a99b05b44aa062ad7df35be9c5839fcab115359ea37e88264fc33994}\\
\code{sw4_infinite_tail_independent_check.py} &
\code{a58da339f5fd7013c130dfaa37d533c5b1263145678cc1152033d9ec892b9ec4}\\
\code{sw4_central_seed_witnesses.csv} &
\code{be4b8cf0fd7bf8c7898beb8698174246d53cf1dd3d4ab64c3d774e1836c6ef3b}\\
\code{sw4_central_seed_capacity.csv} &
\code{04578eb7611359b92a9af9ca6da6a5ed9548a1feaccb2d7724a5e3b0b9903c47}\\
\code{sw4_full_seed_capacity.csv} &
\code{645090a9d9a8b6aab7c03c2116eb1dfb0e53d00b36cf572682b7871738510bc6}\\
\bottomrule
\end{longtable}

\begin{thebibliography}{1}
\bibitem{supplement}
\texttt{[AUTHOR NAME]},
\emph{Supplementary exact certificates for the Steiner--Wiener $4$-index
inverse problem}, accompanying reproducibility archive,
SHA-256
\texttt{1ed15206d460cac319b0b3683a41ade37db687ecd590a199126b08e189b47342}.
\end{thebibliography}

\end{document}
